\chapter{Notation, Conventions, and Status Vocabulary}
\index{notation|textbf}
\index{convention|textbf}

\label{app:notation-and-status}

This appendix is an expository index for the monograph. It does not override the
applicable specifications, proof contracts, software schemas, or retained
evidence. A symbol used by an owning record retains the definition in that
record when it is more specific than the summary below.

\section{Physical and mathematical objects}
\index{physical model}
\index{mathematical operator}
\index{software artifact}


\begin{center}
\small
\begin{tabular}{@{}p{0.20\textwidth}p{0.70\textwidth}@{}}
\toprule
Symbol or term & Meaning in this monograph \\
\midrule
Physical model & A declared material and physical approximation, including its
state variables, branch, and intended interpretation. \\
$\widehat H$ & An abstract Hamiltonian operator on an identified state space;
not a finite matrix unless a representation has been selected. \\
$H$ & A finite matrix representation whose space, basis, geometry, units,
energy reference, spin convention, and gauge must be identified. \\
$\mathcal H$ & An ambient one-particle state space. \\
$\mathcal S$ & A finite retained subspace of an ambient state space. \\
$P_{\mathcal S}$ & Orthogonal projector from the ambient space onto
$\mathcal S$. \\
$\mathcal B$ & An ordered basis for a represented finite space. \\
$N$ or $M$ & A finite represented or retained dimension, as locally defined. \\
$\dagger$ & Adjoint or conjugate transpose in the stated finite-dimensional
inner-product setting. \\
\bottomrule
\end{tabular}
\end{center}

\section{Periodic and Wannier notation}
\index{notation!Bloch and Wannier}


\begin{center}
\small
\begin{tabular}{@{}p{0.20\textwidth}p{0.70\textwidth}@{}}
\toprule
Symbol & Meaning \\
\midrule
$\mathbf r$ & Real-space position coordinate. \\
$\mathbf k$ & Crystal momentum in a declared reciprocal-coordinate convention. \\
$\mathbf R$ & Bravais-lattice translation. \\
$\psi_{n\mathbf k}$ & Bloch state with local band or frame index $n$. \\
$u_{n\mathbf k}$ & Cell-periodic factor of a Bloch state. \\
$U(\mathbf k)$ & Momentum-dependent unitary frame choice within a retained
Bloch subspace; not to be confused with an interspace identification below. \\
$w_{n\mathbf R}$ & Wannier basis function generated from a retained Bloch frame. \\
$H^{(\mathrm W)}(\mathbf k)$ & Hamiltonian matrix in the selected Wannier
frame at $\mathbf k$. \\
$H^{(\mathrm W)}(\mathbf R)$ & Real-space Wannier Hamiltonian block under the
recorded Fourier convention. \\
$N_{\mathbf k}$ & Number of sampled wavevectors in the applicable uniform grid. \\
\bottomrule
\end{tabular}
\end{center}

The same letter $U$ is conventional for both a Wannier gauge and an interspace
identification. The local domain and codomain must disambiguate it. Where
confusion is likely, the monograph uses $G$ for an in-space gauge rotation and
$U$ or $A$ for a map between retained spaces.

\section{Frames, gauges, and alignment}
\index{notation!alignment}


\begin{center}
\small
\begin{tabular}{@{}p{0.20\textwidth}p{0.70\textwidth}@{}}
\toprule
Symbol & Meaning \\
\midrule
$V$ & Rectangular matrix whose orthonormal columns form a retained frame. \\
$G$ & Unitary coordinate change within one retained space. \\
$VV^{\dagger}$ & Ambient-space projector associated with an orthonormal frame. \\
$U$ or $A$ & Explicit identification between two retained spaces; its direction
is stated with its definition. \\
$H_b$ & Pristine or bulk represented Hamiltonian in its declared coordinates. \\
$H_d$ & Doped represented Hamiltonian for dopant $d$. \\
$G_b,G_d$ & Independent unitary coordinate changes in pristine and doped
retained spaces. \\
$\Delta H_{d\to b}$ & Doped operator pulled into pristine coordinates and then
subtracted: $U^{\dagger}H_dU-H_b$. \\
$\Delta H_{b\to d}$ & Pristine operator pushed into doped coordinates and then
subtracted: $H_d-UH_bU^{\dagger}$. \\
$\overline H$ & Operator after application of the explicitly recorded scalar
energy-reference convention. \\
$\theta_j$ & Principal angle between two retained subspaces. \\
\bottomrule
\end{tabular}
\end{center}

A subscript $b$ denotes bulk or pristine material, while $d$ denotes a dopant
branch. When the identity of the dopant matters,
$d\in\{\mathrm P,\mathrm B\}$. Neither notation fixes charge state, spin branch,
or supercell size without an accompanying record.

\section{Reduction notation}

\begin{center}
\small
\begin{tabular}{@{}p{0.20\textwidth}p{0.70\textwidth}@{}}
\toprule
Symbol & Meaning \\
\midrule
$\mathfrak M_j$ & Declared reduced-model class indexed by complexity or another
frozen ordering. \\
$\boldsymbol\theta$ & Parameter vector within a declared model class. \\
$f_{\phi}$ & Learned parameter map with trainable parameters $\phi$ and declared
physical descriptors as inputs. \\
$\mathcal C$ & Deterministic constructor mapping admissible parameters to a
represented Hamiltonian. \\
$\epsilon(\mathfrak M)$ & Best declared normalized discrepancy attainable in
model class $\mathfrak M$; a computed fit need not attain this infimum. \\
$\mathcal L_E$ & Spectral training objective with stated data, weights, units,
and normalization. \\
$\mathcal L_H$ & Aligned represented-operator objective. \\
$\tau_E,\tau_H$ & Prespecified spectral and operator tolerances. \\
$\mathfrak A_{E,j}$ & Spectral-admissible subset of model class $j$. \\
$\mathfrak A_{H,j}$ & Operator-admissible subset of model class $j$. \\
$d_{\mathrm{TB}}$ & Declared metric on compatible canonical tight-binding
coordinates. \\
$\delta_j^{\ast}$ & Minimum or infimum separation between the two admissible
sets, according to the established assumptions. \\
$\mathcal R$ & Reduction map with an explicitly stated domain and codomain. \\
$\varepsilon_{\mathrm{path}}$ & Norm of the difference between two aligned
reduction paths. \\
\bottomrule
\end{tabular}
\end{center}

The words \emph{projection}, \emph{disentanglement}, \emph{basis
transformation}, \emph{localization}, \emph{truncation}, \emph{fitting}, and
\emph{reduction} denote different operations. They are not used as synonyms.

\section{Excluded-space and continuum notation}

\begin{center}
\small
\begin{tabular}{@{}p{0.20\textwidth}p{0.70\textwidth}@{}}
\toprule
Symbol & Meaning \\
\midrule
$P,Q$ & Complementary orthogonal projectors with $P+Q=I$. A superscript $(P)$
marks restriction to the retained $P$ space; phosphorus is written $\mathrm P$
when ambiguity is possible. \\
$H_{PP},H_{PQ}$ & Blocks such as $PHP$ and $PHQ$ in a declared decomposition. \\
$H_{\mathrm{eff}}(E)$ & Energy-dependent Feshbach effective operator in the
retained $P$ space. \\
$\Delta_Q$ & Distance from $E$ to the spectrum of $H_{QQ}$ in the declared
self-adjoint resolvent setting. \\
$V_{\mathrm{cont},d}$ & Candidate continuum impurity operator for dopant $d$. \\
$D_d$ & Aligned atomistic-minus-continuum discrepancy. \\
$P_{>R}$ & Projector onto a declared exterior region beyond radius $R$. \\
$\eta_d(R)$ & Exterior residual $\|P_{>R}D_dP_{>R}\|$ in a declared norm. \\
$r_{c,d}(\tau)$ & Tolerance-dependent crossover infimum; not a universal
material constant. \\
$J$ & Continuum-to-atomistic or continuum-to-Wannier embedding, with its
normalization and validation domain stated separately. \\
\bottomrule
\end{tabular}
\end{center}

\citationtodo{low priority}{Check the naming and definition of the
energy-dependent Feshbach effective operator against Feshbach (1958),
bibliography key \texttt{feshbach1958} (already present).}

\section{Norms and residuals}
\index{norm}
\index{residual}


\begin{align}
  \|A\|_{\mathrm F}
  &=
  \left(\sum_{i,j}|A_{ij}|^2\right)^{1/2},\\
  \|A\|_2
  &=
  \sup_{x\ne0}\frac{\|Ax\|_2}{\|x\|_2},\\
  \|A\|_{\max}
  &=
  \max_{i,j}|A_{ij}|.
\end{align}

The Frobenius and induced Euclidean operator norms are invariant under unitary
conjugation. Entrywise maxima, shell decompositions, and spatial truncations may
have narrower invariance groups. Every reported residual states whether it is
absolute or normalized and gives its units and zero-scale behavior.

\section{Proof status}

\begin{description}
  \item[\texttt{proposed}.] Statement, definitions, or strategy exists without a
    complete established proof.
  \item[\texttt{in development}.] Proof construction is active and immediate
    prerequisites are sufficiently identified.
  \item[\texttt{blocked}.] A named missing assumption, definition, prerequisite,
    or human scientific decision prevents progress.
  \item[\texttt{complete, unreviewed}.] A complete derivation is recorded but
    has not passed independent mathematical review.
  \item[\texttt{reviewed}.] The derivation and declared prerequisites passed
    independent mathematical review.
  \item[\texttt{numerically checked}.] Separate numerical evidence agrees under
    declared conditions; this supplements rather than replaces proof.
  \item[\texttt{retired}.] The claim was superseded, merged, or rejected with a
    recorded disposition.
\end{description}

A mechanization contract may separately be \texttt{frozen}, \texttt{encoded},
\texttt{checked}, or \texttt{cross-checked} in each backend. These words do not
replace the proof-package status. In particular, the currently checked Lean
encoding of \texttt{PRF-05.01} does not make the entire \texttt{PRF-05} package
complete or reviewed.

\section{Scientific and evidentiary status}
\index{evidence!status vocabulary}


\begin{description}
  \item[Calculated result.] Produced by an identified calculation with retained
    provenance.
  \item[Literature value.] Taken from an identified and verified source.
  \item[Expected behavior.] A prediction, not an observation.
  \item[Illustrative example.] Explanatory material without research-evidence
    status.
  \item[Synthetic test data.] Constructed for software or numerical testing.
  \item[Placeholder.] Incomplete material awaiting authoritative content.
  \item[Proposed work.] Planned but not completed.
\end{description}

Mathematical proof, software verification, numerical verification, scientific
validation, and uncertainty quantification remain distinct evidence types. A
claim may require more than one of them, and none is inferred merely from the
presence of the others.

\section{Units and coordinate conventions}
\index{units}
\index{coordinate convention}


Every physical numerical value states its unit. Energy references used for
plots, bulk observables, and pristine--doped subtraction are identified
separately. Real-space coordinates identify lattice vectors, cell convention,
and site mapping. Reciprocal coordinates identify the reciprocal basis and any
path parameterization. Spinless, collinear-spin, and noncollinear spinor
representations are never mixed implicitly.

The notation in this appendix is intentionally insufficient to authorize a
comparison. The represented records and owning specifications must still supply
the exact metadata needed by the operation.
